% =============================================================================
% DISCUSSION — 9-condition study, populated from the result bundle.
% Findings (5.1) and the cost-model/tool-forcing interpretation (5.2-5.3)
% draw on results/analysis/*; threats (5.5) and scope (5.7) are general.
% =============================================================================

\subsection{Principal Findings}
\label{sec:disc_summary}

We restate the four pre-registered hypotheses
(\cref{sec:intro_hypothesis}) and the outcome of each, drawing only on
the contrasts in \cref{sec:results_primary,sec:results_architecture}.

\begin{itemize}\setlength\itemsep{3pt}
  \item \textbf{H1/H2 (does best-practice verifier-guided extraction
        beat B$'$?).} No. Both upfront ($72.5\%$ vs $71.6\%$;
        $p_{\text{Bonf}} = 1.0$) and reactive ($72.2\%$;
        $p_{\text{Bonf}} = 1.0$) are statistically indistinguishable
        from B$'$. Neither cleared the Bonferroni threshold; neither
        came close.
  \item \textbf{H3 (citations).} No. Citation grounding does not move
        accuracy over reactive ($+0.3$ pp, $p = 0.80$) --- and half of
        all source spans failed substring validation, so it mostly
        removed fields from the comparison.
  \item \textbf{H4 ($k$-shot).} No. $k=3$ voting is if anything
        slightly worse than single-sample reactive ($-0.2$ pp,
        $p = 0.89$), despite the samples genuinely diverging
        ($23.7\%$ of field-votes reached no majority).
  \item \textbf{Anchor (does tool access help?).} Yes, decisively, and
        this is where the accuracy lives: tool access (B vs A,
        $+8.7$ pp) and especially forced tool use (B$'$ vs B,
        $+21.3$ pp) are the only significant primary contrasts. They
        bound the room a verifier has to improve and are not themselves
        claims about verification.
\end{itemize}

The four primary hypotheses therefore all fail to reject the null: at
best-practice configuration the input verifier sits on top of B$'$
without moving accuracy. The rest of this section explains \emph{why}
the gap is null --- and why that is the expected outcome once the
reference layer is taken seriously.

\subsection{Interpreting a Residual Gap: The Reference Layer}
\label{sec:disc_extractor}

The verifier in this study is deterministic
(\cref{sec:methods_verifier}): exact match on enums, numeric tolerance
on numbers, canonical-form match on booleans, and no flag at all when
the extractor abstains. None of that machinery can synthesize a wrong
flag on its own --- a spurious flag requires a wrong reference, and the
reference comes from the LLM extractor. The extracted reference is
therefore the load-bearing component, and any gap between a
best-practice extractor pipeline and B$'$ localizes there rather than
in the comparison logic.

We state the limit as a cost model. Where the verifier's reference is
produced by an extractor $M_e$ over a required-field set $\mathcal{F}$,
the verifier's net effect on accuracy is governed by the joint
distribution of three terms:
(a) the primary model's input-error rate on $\mathcal{F}$ --- the
errors a true flag can fix;
(b) $M_e$'s extraction-error rate on $\mathcal{F}$ --- the errors that
produce false flags;
(c) the asymmetric cost of a false flag (which re-prompts and can flip
a correct call to wrong) versus a true flag (which can fix a wrong
call).
Verification helps only when true flags outweigh false flags under
this cost. The schema gate (\cref{sec:methods_verifier}) restricts the
comparison to required fields, which shrinks the surface on which term
(b) can act; the four extractor pipelines vary placement, output
format, and sampling, each of which is a different attempt to lower
term (b). Empirically, none of the three levers moved accuracy, so none
measurably shifted the balance of terms (a)--(c). The flag statistics
are consistent with this: reactive extraction fired \emph{more} often
than upfront ($16.1\%$ vs $12.2\%$ of vignettes) yet had a \emph{lower}
correct-on-intervention rate ($66.9\%$ vs $74.6\%$), so its extra
interventions were not net-corrective. The likeliest reading is that,
once tool use is forced, the primary model's input-error rate on
required fields (term a) is already low --- B$'$ alone reaches
$71.6\%$ --- so there are few wrong calls for a true flag to rescue,
while the residual false flags from term (b) roughly cancel whatever
true flags remain. The verifier lands at parity because, post-gating,
both kinds of flag are rare and approximately balanced.

A subtlety the cost model alone misses: ``extraction error'' is
broader than ``the extractor read the wrong value.'' A flag can fire
when the extractor's value is semantically correct but serialized
differently from the model's (the \texttt{sex=female} vs
\texttt{sex=1} pattern traced in \cref{app:loop_trace}). The binding
constraint there is the join of (extractor canonical form, verifier
comparison strictness, model serialization) --- any one of the three
can drive a flag the other two cannot resolve. This residue persists
under schema-gating: \texttt{sex} is still the most-flagged field
($22.6\%$ of upfront flags, co-top under reactive) --- the same field
that drives the \texttt{bmi\_001} trace, where the extractor's
semantically-correct \texttt{female} and the model's \texttt{1} disagree
only in encoding. Schema-gating removed flags on fields the calculator
ignores, but it cannot reconcile encoding mismatches on \emph{required}
fields, so this join keeps term (b)'s false-flag floor above zero and
is one concrete reason the verifier cannot get ahead of B$'$.

\subsection{Why Tool-Forcing Amplifies the Reference}
\label{sec:disc_toolforcing}

The verifier fires once per tool call, so its leverage scales with the
tool-call rate. The B$'$ prompt forces a tool call on every vignette
and removes the model's option to decline when uncertain; on every
such call the verifier injects feedback carrying the extracted
reference. This amplifies whatever the reference does: a reliable
reference is applied more often (more fixes), and an unreliable one is
applied more often (more mis-corrections). The amplification is
multiplicative in the tool-call rate, and its sign is set entirely by
the reference quality of \cref{sec:disc_extractor} --- forcing is not
harmful or helpful on its own. Here the amplification ran to a wash:
forcing raised the verifier's exposure as designed (the model emits
$\sim$$3.7$ tool calls per vignette and the verifier intervened on
$12.2\%$--$16.1\%$ of vignettes), but with post-gating flags roughly
balanced between true and false, the larger exposure moved net accuracy
in neither direction. This is the benign face of the \emph{same}
amplification that, un-gated, produced the pilot's $\sim$$9$ pp harm
(\cref{sec:methods_pilots}): schema-gating shrank the false-flag surface
enough that forcing no longer multiplies a net-negative signal.

\subsection{Implications for Deploying Verifier-Guided Reasoning in Clinical AI}
\label{sec:disc_implications}

Two recommendations follow from the mechanism regardless of the
headline numbers; a third is conditional on them.

\textit{Validate the reference pipeline independently before adopting
intermediate verification.} The verifier's correctness guarantee is no
stronger than the extractor's field-level accuracy on the deployment's
live field distribution. A held-out, hand-annotated field-accuracy
gate on the fields the deployment actually uses --- not the union of
all possible MCP fields --- is the precondition the cost model in
\cref{sec:disc_extractor} implies.

\textit{Prefer a non-LLM reference where one exists.} Where structured
patient data are available (EHR fields, device vitals, LIS results),
they remove term (b) for the fields they carry. Reserve LLM extraction
for fields the structured record does not hold.

\textit{Choose the extractor pipeline on the measured cost--accuracy
frontier.} Among the verifier pipelines, reactive extraction dominates:
it matches upfront's accuracy at $\sim$$40\%$ of the token cost
($32$K vs $79$K per vignette) and is faster than $k$-shot
(\cref{sec:results_cost}). But on this benchmark \emph{no} pipeline is
Pareto-efficient against B$'$ itself, which matches their accuracy at
lower cost --- so input verification earns its overhead only where the
reference can be made more reliable than a single best-practice LLM
extractor, i.e., under the first two recommendations above.

\subsection{Threats to Validity}
\label{sec:disc_limitations}

\textit{Single primary model.} All conditions evaluate
Qwen3-30B-A3B-Thinking-2507. The mechanism we identify (the extracted
reference as the binding component) should transfer across models, but
the magnitudes will not.

\textit{Single benchmark.} medical\_calculator\_eval covers
deterministic calculator dispatch with typed inputs. Generalization to
differential-diagnosis, triage, or open-ended tool surfaces --- where
inputs are less structured --- is untested.

\textit{Single extractor model.} Sonnet 4.6 is one extractor; another
could show different per-field error rates and thus a different verdict.
The claim is that \emph{some} LLM extractor is the binding constraint,
not that Sonnet specifically fails.

\textit{No ground-truth-reference comparator.} No condition consults a
hand-annotated reference, which would isolate the verifier mechanism
from extractor fallibility entirely. This is the most important
comparator the study lacks (\cref{sec:disc_followups}).

\textit{Conservative verifier is silent on extractor misses.} The ``no
extractor value $\rightarrow$ no flag'' rule (\cref{sec:methods_verifier})
suppresses false positives but makes a silent extractor (one that
extracts nothing for a field the model proposes) invisible to the
firing-rate statistics. Some fraction of vignettes may carry silent
verifier failures the firing counts cannot see.

\textit{Extractor stochasticity.} The extractor runs at temperature
$0.7$ (and $k=3$ in one pipeline), so its reference is not
deterministic across runs. The $k$-shot diagnostics quantify the
variance directly: across $20{,}271$ field-votes, $61.1\%$ showed full
three-way agreement, $15.2\%$ a $2/3$ majority, and $23.7\%$ no majority
at all. So the reference does vary across samples on a meaningful
minority of fields --- but that variance falls on ambiguous fields,
not on the common ones (\texttt{sex}, \texttt{age}, \texttt{height})
that dominate both the flag counts and the accuracy-relevant decisions.
A repeated run would shift total flag counts slightly without changing
the architectural pattern or the null verdict.

\textit{Apparatus-gate scope.} The gate (\cref{sec:methods_gate})
detects drift on the harness, dataset, prompt rendering, and MCP
transport via Sonnet on Condition B. It does not independently validate
the primary Qwen3 numbers; we assume the harness behaves identically
for the primary model conditional on the gate holding.

\subsection{Future Directions}
\label{sec:disc_followups}

The study fixes the verifier and varies only the extractor pipeline.
Three reference-layer variants are the natural next experiments, each
holding all other variables at this study's values. We list them as
hypothesized future work, not as findings.

\begin{itemize}\setlength\itemsep{2pt}
  \item \textbf{Hand-annotated ground-truth reference.} Removes the LLM
        extractor entirely and answers whether procedural input
        verification helps when the reference is reliable. Cost:
        annotation of all $1{,}066$ vignettes.
  \item \textbf{Multi-LLM ensemble extractor.} Requires agreement
        across $\ge 2$ extractor models before the verifier's reference
        is populated, targeting single-model variance in term (b).
  \item \textbf{Structured-record reference.} A non-LLM reference layer
        (e.g., EHR-keyed inputs) for a deployment-realistic setting;
        requires a benchmark that ships structured records alongside
        vignettes, which medical\_calculator\_eval does not.
\end{itemize}

\subsection{What This Paper Does Not Claim}
\label{sec:disc_notclaim}

To preempt common misreadings, we name the claims this paper does
\emph{not} make.

\textbf{1. ``Interwhen does not work.''} Interwhen's verify/fix
mechanism operates as designed in every condition we run. Our
adaptation repurposes its LLM extraction step to recover an independent
reference from the case rather than the model's asserted state from the
trace (\cref{sec:bg_interwhen}); any limit we find is in that reference
layer, not the mechanism.

\textbf{2. ``Output verification is unhelpful.''} We study procedural
\emph{input} verification specifically (\cref{sec:bg_surfaces}). Output
verification, outcome verification, grounded retrieval, and
human-in-the-loop regimes are different surfaces; Condition D is the
only output-style comparator here and is scoped as a control, not the
object of study.

\textbf{3. ``Sonnet is a poor extractor.''} Sonnet 4.6 is near the
frontier of clinical-text extraction. Any underperformance we attribute
to the extractor is a property of the prompt regime (e.g., a
$\sim$500-field upfront schema), not the model's underlying capability.

\textbf{4. ``Tool-forcing prompts are dangerous.''} Forcing amplifies
whatever follows it in both directions (\cref{sec:disc_toolforcing}).
With a reliable reference the amplification is beneficial; the prompt
is not harmful on its own.

\textbf{5. ``Any one pipeline is the right deployment configuration.''}
Each contrast is one mechanistic test. A production choice depends on
the measured cost--latency frontier, field-level validation of the
extractor, and the specific clinical workflow --- more than a single
accuracy contrast can settle.
